\documentclass[a4paper]{arrowhead}

\usepackage[yyyymmdd]{datetime}
\usepackage{etoolbox}
\usepackage[T1]{fontenc}
\usepackage[utf8]{inputenc}
\usepackage{multirow}
\usepackage{longtable}
\usepackage{amssymb}
\usepackage{float}
\usepackage{tikz}
\usetikzlibrary{positioning,arrows.meta,fit,calc,shapes.misc,shapes.geometric}

\renewcommand{\dateseparator}{-}
\setlength{\emergencystretch}{3em}
\lstset{basicstyle=\ttfamily\footnotesize,breaklines=true,columns=fullflexible,frame=single,rulecolor=\color{gray!50},backgroundcolor=\color{gray!6}}
\tikzset{
  ahsys/.style={draw=ArrowheadBlue,very thick,fill=ArrowheadBlue!6,rounded corners=2pt,minimum width=3.2cm,minimum height=1.1cm,align=center,font=\small\bfseries},
  ahext/.style={draw=ArrowheadGrey,thick,dashed,fill=gray!5,rounded corners=2pt,minimum width=3.2cm,minimum height=1.0cm,align=center,font=\small},
  ahsrv/.style={draw=ArrowheadPurple,thick,fill=ArrowheadPurple!8,rounded corners=6pt,minimum height=0.55cm,inner xsep=4pt,align=center,font=\scriptsize\ttfamily},
  ahcsm/.style={draw=ArrowheadTeal,thick,fill=ArrowheadTeal!8,rounded corners=6pt,minimum height=0.55cm,inner xsep=4pt,align=center,font=\scriptsize\ttfamily},
  ahdb/.style={draw=ArrowheadOrange,thick,fill=ArrowheadOrange!8,cylinder,shape border rotate=90,aspect=0.25,minimum height=0.9cm,align=center,font=\scriptsize},
  ahflow/.style={-{Stealth[length=2.2mm]},thick,ArrowheadGrey},
  ahlabel/.style={font=\scriptsize,fill=white,inner sep=1pt,align=center},
}

%% Special references
\newcommand{\mref}[1]{{\textcolor{ArrowheadPurple}{\hyperref[sec:model:#1]{#1}}}}
\newcommand{\oref}[1]{{\textcolor{ArrowheadBlue}{\hyperref[sec:operations:#1]{#1}}}}
\newcommand{\pdef}[1]{{\textcolor{ArrowheadGrey}{#1\label{sec:model:primitives:#1}\label{sec:model:primitives:#1s}\label{sec:model:primitives:#1es}}}}
\newcommand{\pref}[1]{{\textcolor{ArrowheadGrey}{\hyperref[sec:model:primitives:#1]{#1}}}}

\newrobustcmd\osubsection[3]{
  \addtocounter{subsection}{1}
  \addcontentsline{toc}{subsection}{\protect\numberline{\thesubsection}operation \textcolor{ArrowheadBlue}{#1}}
  \subsection*{
    \thesubsection\quad
    operation
    \textcolor{ArrowheadBlue}{#1}
    (\notblank{#2}{\mref{#2}}{})
    \notblank{#3}{: \mref{#3}}{}
  }
  \label{sec:operations:#1}
}
\newrobustcmd\msubsection[2]{
  \addtocounter{subsection}{1}
  \addcontentsline{toc}{subsection}{\protect\numberline{\thesubsection}#1 \textcolor{ArrowheadPurple}{#2}}
  \subsection*{\thesubsection\quad#1 \textcolor{ArrowheadPurple}{#2}}
  \label{sec:model:#2} \label{sec:model:#2s} \label{sec:model:#2es}
}

\begin{document}

%% Arrowhead Document Properties
\ArrowheadTitle{workflowRecipe}
\ArrowheadServiceID{workflowRecipe}
\ArrowheadType{Service Description}
\ArrowheadTypeShort{SD}
\ArrowheadVersion{5.4.0}
\ArrowheadDate{2026-11-23}
\ArrowheadAuthor{Documentation maintenance}
\ArrowheadStatus{REVIEW DRAFT}
\ArrowheadContact{Repository maintainers}
\ArrowheadFooter{\href{www.arrowhead.eu}{www.arrowhead.eu}}
\ArrowheadSetup
%%

%% Front Page
\begin{center}
  \vspace*{1cm}
  \huge{\arrowtitle}

  \vspace*{0.2cm}
  \LARGE{\arrowtype}
  \vspace*{1cm}

  \Large{Service ID: \textit{``\arrowid''}}
  \vspace*{\fill}

  

  \vspace*{1cm}
  \vspace*{\fill}

  % Front Page Abstract
  \begin{abstract}
    This document describes the \texttt{workflowRecipe} service of the Eclipse Arrowhead workflow support systems, which validates, registers, lists, retires and compiles workflow recipes. It is provided by the WorkflowManager system, service version 1.0.0, documentation edition 5.4.0.
  \end{abstract}

  \vspace*{1cm}
\end{center}

\newpage
%%

%% Table of Contents
\tableofcontents
\newpage
%%

\section{Overview}
\label{sec:overview}
The \texttt{workflowRecipe} service is the recipe catalogue of a WorkflowManager. A recipe is the
design-time description of a workflow: which Arrowhead services are consumed (service definitions and capabilities,
never endpoints), in which order and under which conditions, and how data flows between the exchanges. The service
validates recipes against the well-formedness conditions of the formal recipe model~\cite{varga2026} and the
WorkflowManager's own checks, stores accepted recipes immutably by name and version, and compiles them to plans of the
existing Choreographer.

The rest of this document is organized as follows.
In Section~\ref{sec:operations}, we describe the abstract message operations provided by the service.
In Section~\ref{sec:model}, we end the document by presenting the data types used by the mentioned operations.

\subsection{Significant Prior Art}
The recipe model is the formal definition $W = (P, A, S, \theta, \iota, \mu, q, pre, \eta, qnty, dep, \kappa)$
of Varga~\cite{varga2026}, with its well-formedness conditions (W1)--(W4). Separating the workflow description from the
consumers that execute it follows Derhamy's design principles of loose coupling and specialisation~\cite{derhamy2018}
and the workflow management system of Garcia Represa~\cite{represa2022}. The plan compiled by \oref{compile} is the
plan model of the Arrowhead Choreographer (4.6.2 to 5.4.0).

\subsection{How This Service Is Meant to Be Used}
An engineering tool, normally the WorkflowConfigurator, calls \oref{validate} while a workflow is being
designed and \oref{register} when it is accepted. A product, line controller or upper WorkflowManager then starts
sessions of the registered recipe through \texttt{workflowInstance}. Old versions are retired with \oref{retire}: they
remain readable for traceability but cannot be started. A recipe that should run on an existing Choreographer is
compiled with \oref{compile} and submitted to the Choreographer's plan management.

\subsection{Important Delimitations and Dependencies}
\label{sec:delimitations}
\begin{itemize}
\item Recipes are immutable: registering different content under an existing name and version is refused (409). A change needs a new version.
\item The validator decides on the structure, the data flow, the lease discipline and, in a connected deployment, on the availability of services and executors in the running cloud (checks C01, C02). It does not check the semantic correctness of the equipment operations.
\item \texttt{quantity} $>$ 1 is not accepted (M01); repeated execution is outside the proven model.
\item Compilation to a Choreographer plan is lossy: inputs ($\kappa_s$), guards, dependencies and policies cannot be expressed in a plan and are listed in \texttt{lossySteps}.
\item Retiring is the only life-cycle change; recipes are never deleted through this service.
\end{itemize}

\paragraph{Authorization.} Each operation is authorized separately (\texttt{workflowRecipe.validate}, \ldots). The WorkflowConfigurator needs \texttt{validate}, \texttt{register}, \texttt{query} and \texttt{get}.

\newpage
\section{Service Interface}
\label{sec:operations}
This section describes the interface of the \texttt{workflowRecipe} service. Each subsection names an abstract operation,
an input type and an output type, in that order. The input type is named inside parentheses, while the output type is
preceded by a colon. Input and output types are only denoted when accepted or returned, respectively, by the operation
in question. All abstract data types named in this section are defined in Section~\ref{sec:model}. The concrete
interface (HTTP/TLS/JSON) is specified in the IDD of the same name.

\begin{figure}[ht!]
  \centering
  \resizebox{0.92\textwidth}{!}{
\begin{tikzpicture}
\node[ahsys] (prov) at (4.6,0) {WorkflowManager\\[-1pt]{\normalfont\scriptsize $\ll$provider$\gg$}};
\node[ahsrv] (srv) at (1.6,0) {workflowRecipe};
\draw[thick,ArrowheadPurple] (srv.east) -- (prov.west);
\node[ahext] (c0) at (-3.0,1.20) {WorkflowConfigurator\\[-1pt]{\scriptsize $\ll$consumer$\gg$}};
\draw[ahflow] (c0.east) -- (srv.west);
\node[ahext] (c1) at (-3.0,0.00) {Engineering tool\\[-1pt]{\scriptsize $\ll$consumer$\gg$}};
\draw[ahflow] (c1.east) -- (srv.west);
\node[ahext] (c2) at (-3.0,-1.20) {Upper WorkflowManager\\[-1pt]{\scriptsize $\ll$consumer$\gg$}};
\draw[ahflow] (c2.east) -- (srv.west);
\end{tikzpicture}
  }
  \caption{SysML block view of the \texttt{workflowRecipe} service: provider and typical consumers.}
  \label{fig:SD}
\end{figure}

The following interface operations are available.

{\small
\begin{longtable}{|p{4.4cm}|p{5.6cm}|p{5.8cm}|}

\hline
\rowcolor{gray!33} Operation & Input & Output \\ \hline
\endhead
\oref{validate} & \mref{Recipe} & \mref{ValidationReport} \\ \hline
\oref{register} & \mref{Recipe} & \mref{RecipeRegistration} \\ \hline
\oref{query} & \mref{RecipeFilter} & \mref{RecipeList} \\ \hline
\oref{get} & \mref{RecipeKey} & \mref{RecipeRecord} \\ \hline
\oref{retire} & \mref{RecipeKey} & \mref{RecipeRecord} \\ \hline
\oref{compile} & \mref{RecipeKey} & \mref{ChoreographerPlan} \\ \hline
\end{longtable}
}

\osubsection{validate}{Recipe}{ValidationReport}
Validates a recipe without storing it. Every finding is reported; the verdict is the worst outcome.
\osubsection{register}{Recipe}{RecipeRegistration}
Validates the recipe and stores it when the verdict is ACCEPT. The registering system is recorded as \texttt{registeredBy}.
\osubsection{query}{RecipeFilter}{RecipeList}
Lists registered recipes in registration order, optionally only the versions of one recipe.
\osubsection{get}{RecipeKey}{RecipeRecord}
Returns one registered recipe version, active or retired.
\osubsection{retire}{RecipeKey}{RecipeRecord}
Marks a recipe version \texttt{RETIRED}. Running sessions continue; new sessions are refused.
\osubsection{compile}{RecipeKey}{ChoreographerPlan}
Compiles an active recipe version (the latest when \texttt{version} is absent) to a Choreographer plan request.

\subsection{Guard expressions}
\label{sec:guards}
A step may put a condition on the transition to each successor (\texttt{guards}, extension x-guard). The language is
deliberately small and side-effect free, so a WorkflowManager can evaluate it from $\rho$ and $\xi$ alone and all
implementations can share the grammar:
\begin{lstlisting}
expr      := or
or        := and ( "||" and )*
and       := unary ( "&&" unary )*
unary     := "!" unary | compare
compare   := operand ( ( "==" | "!=" | "<" | "<=" | ">" | ">=" ) operand )?
operand   := "(" expr ")" | reference | number | string | "true" | "false" | "null"
reference := "${param." name path "}" | "${result." step path "}"   path := ( "." segment )*
string    := '...' | "..."   (a backslash escapes the next character)
\end{lstlisting}
A lone operand is true when it is Boolean true, a non-zero number, or a non-empty string, array or object. A missing
reference is null. Numbers compare numerically, strings lexicographically, and \texttt{==}/\texttt{!=} compare JSON
values. Ordering a number with a non-number is false.

A step is enabled when every predecessor is resolved (completed or skipped) and at least one transition into it was
taken; it is skipped (x-skip) when none was taken. Without guards this is exactly the rule S-Enable of the model.


\subsection{Validator checks}
\label{sec:codes}
The validator applies the checks below and reports each violation as a \mref{Finding}. The S, D and M checks realise the
well-formedness conditions (W1)--(W4) of~\cite{varga2026}; the L checks the premise of Proposition~1 in its lease
variant; the G checks the guarded-transition extension; C01 and C02 run only in a WorkflowManager connected to a
ServiceRegistry. The decision rule is REJECT $>$ CLARIFY $>$ REVISE $>$ ACCEPT, and only ACCEPT can be registered.

{\small
\begin{longtable}{|p{1.3cm}|p{1.9cm}|p{12.8cm}|}

\hline
\rowcolor{gray!33} Code & Outcome & Meaning \\ \hline
\endhead
\texttt{S00} & REVISE & Name or version not valid. \\ \hline
\texttt{S01} & REVISE & Two actions have the same name. \\ \hline
\texttt{S02} & REVISE & Two steps have the same name. \\ \hline
\texttt{S03} & REVISE & The first action does not exist. \\ \hline
\texttt{S04} & REVISE & The next action does not exist. \\ \hline
\texttt{S05} & REVISE & An action cannot be reached, or actions form a loop. \\ \hline
\texttt{S06} & REVISE & Transition to a step in another action. \\ \hline
\texttt{S07} & REVISE & The steps form a loop. \\ \hline
\texttt{S08} & REVISE & The workflow or an action is empty. \\ \hline
\texttt{S09} & REVISE & A step has no service. \\ \hline
\texttt{P01} & REVISE & Two session parameters have the same name. \\ \hline
\texttt{D01} & REVISE & Dependency on an unknown step. \\ \hline
\texttt{D02} & REVISE & A dependency does not run first. \\ \hline
\texttt{D03} & REVISE & The input reads a result the step cannot see (outside $dep(s)$). \\ \hline
\texttt{D04} & REVISE & The input needs a service result no dependency provides (\texttt{requiresResults}). \\ \hline
\texttt{D05} & CLARIFY & The input reads an undeclared session parameter. \\ \hline
\texttt{G01} & REVISE & Condition on a transition that does not exist. \\ \hline
\texttt{G02} & REVISE & A condition cannot be parsed. \\ \hline
\texttt{G03} & REVISE & A condition reads a result that may not exist yet. \\ \hline
\texttt{G04} & ACCEPT & Information: the workflow uses conditions (x-guard). \\ \hline
\texttt{G05} & CLARIFY & A condition reads an undeclared session parameter. \\ \hline
\texttt{L01} & REVISE & Lease of an undeclared resource. \\ \hline
\texttt{L02} & REVISE & No reserve step before the tool use. \\ \hline
\texttt{L03} & REVISE & A tool use does not depend on its reserve step. \\ \hline
\texttt{L04} & CLARIFY & Lease duration missing or shorter than a tool step. \\ \hline
\texttt{L05} & REVISE & No release step after the tool use. \\ \hline
\texttt{M01} & REVISE & Quantity above 1. \\ \hline
\texttt{M02} & REVISE & Executor choice without a capability. \\ \hline
\texttt{C01} & CLARIFY & The service is not found in the running cloud (WorkflowManager catalogue check). \\ \hline
\texttt{C02} & CLARIFY & No executor can run the step (WorkflowManager catalogue check). \\ \hline
\end{longtable}
}


\newpage
\section{Information Model}
\label{sec:model}
Here, all data objects that can be part of the \texttt{workflowRecipe} service are listed in alphabetic order. Each
subsection, which describes one type of object, begins with the \textit{struct} keyword (a collection of named fields,
each with its own data type), the \textit{enum} keyword (a closed set of values) or the \textit{alias} keyword. As a
complement to the explicitly defined types in this section, there is also a list of implicit primitive types in
Section~\ref{sec:model:primitives}.

\msubsection{struct}{Action}
An action $\alpha$: a group of steps run in one work centre.

{\small
\begin{longtable}{|p{3.9cm}|p{4.4cm}|p{7.5cm}|}

\hline
\rowcolor{gray!33} Field & Type & Description \\ \hline
\endhead
\texttt{name} & \pref{Name} & Action name. \\ \hline
\texttt{next} & \pref{Name} & The action that becomes active when this one finishes; absent for the last action. Optional. \\ \hline
\texttt{workCentre} & \pref{String} & Informative work-centre identifier. Optional. \\ \hline
\texttt{dispatch} & \mref{DispatchMode} & Dispatch mode of this action's steps. Optional. \\ \hline
\texttt{steps} & \pref{List}$<$\mref{Step}$>$ & The steps $S_\alpha$. \\ \hline
\end{longtable}
}

\msubsection{struct}{ChoreographerPlan}
A recipe compiled to the Choreographer plan request (\texttt{PlanRequestDTO}), for option A of the coexistence with the Choreographer.

{\small
\begin{longtable}{|p{3.9cm}|p{4.4cm}|p{7.5cm}|}

\hline
\rowcolor{gray!33} Field & Type & Description \\ \hline
\endhead
\texttt{name} & \pref{Name} & Plan name: the recipe name. \\ \hline
\texttt{firstActionName} & \pref{Name} & First action. \\ \hline
\texttt{actions} & \pref{List}$<$\mref{PlanAction}$>$ & Actions. \\ \hline
\texttt{lossySteps} & \pref{List}$<$\pref{String}$>$ & Informational: steps whose input, guards, dependencies or policy the plan cannot express. Not part of the Choreographer request. Optional. \\ \hline
\end{longtable}
}

\msubsection{enum}{DispatchMode}
How a step reaches an executor.

{\small
\begin{longtable}{|p{3.9cm}|p{11.9cm}|}

\hline
\rowcolor{gray!33} Value & Meaning \\ \hline
\endhead
\texttt{push} & The WorkflowManager selects an executor and calls \texttt{workflowExecution.start} (default). \\ \hline
\texttt{pull} & The WorkflowManager publishes the work item; an executor claims it with \texttt{workflowWorkItem.claim}. \\ \hline
\texttt{event} & Reserved for EventHandler notification; currently handled as \texttt{pull}. \\ \hline
\end{longtable}
}

\msubsection{struct}{ExecutorRequirement}
$\eta(s)$: constraints on the executor that runs a step.

{\small
\begin{longtable}{|p{3.9cm}|p{4.4cm}|p{7.5cm}|}

\hline
\rowcolor{gray!33} Field & Type & Description \\ \hline
\endhead
\texttt{capability} & \pref{String} & A capability the executor must declare in \texttt{workflowCapabilities}. Optional. \\ \hline
\texttt{workstation} & \pref{String} & The executor's workstation. Optional. \\ \hline
\texttt{system} & \pref{SystemName} & A named executor system. Optional. \\ \hline
\end{longtable}
}

\msubsection{struct}{Finding}
One validator finding.

{\small
\begin{longtable}{|p{3.9cm}|p{4.4cm}|p{7.5cm}|}

\hline
\rowcolor{gray!33} Field & Type & Description \\ \hline
\endhead
\texttt{code} & \pref{String} & Check code, e.g. \texttt{D03}; the codes are listed in the \texttt{workflowRecipe} SD. \\ \hline
\texttt{outcome} & \mref{Verdict} & Outcome of this finding. \\ \hline
\texttt{basis} & \pref{String} & The rule the check enforces, e.g. \texttt{W1}, \texttt{V9}, \texttt{schema}. \\ \hline
\texttt{location} & \pref{String} & Where the finding applies, e.g. \texttt{step evaluate}. \\ \hline
\texttt{message} & \pref{String} & Human-readable explanation. \\ \hline
\end{longtable}
}

\msubsection{enum}{LeaseMode}
Role of a step for a shared resource.

{\small
\begin{longtable}{|p{3.9cm}|p{11.9cm}|}

\hline
\rowcolor{gray!33} Value & Meaning \\ \hline
\endhead
\texttt{reserve} & Reserves the resource; its result carries the lease token. \\ \hline
\texttt{use} & Uses the resource; the executor renews the lease immediately before the invocation. \\ \hline
\texttt{release} & Releases the resource. \\ \hline
\end{longtable}
}

\msubsection{struct}{LeasePolicy}
The role of a step for a shared resource.

{\small
\begin{longtable}{|p{3.9cm}|p{4.4cm}|p{7.5cm}|}

\hline
\rowcolor{gray!33} Field & Type & Description \\ \hline
\endhead
\texttt{resource} & \pref{Name} & A declared resource. \\ \hline
\texttt{mode} & \mref{LeaseMode} & \texttt{reserve}, \texttt{use} or \texttt{release}. \\ \hline
\end{longtable}
}

\msubsection{enum}{OnError}
Reaction to an \texttt{ERROR} report.

{\small
\begin{longtable}{|p{3.9cm}|p{11.9cm}|}

\hline
\rowcolor{gray!33} Value & Meaning \\ \hline
\endhead
\texttt{rebind} & Exclude the executor for this session and bind the step again. \\ \hline
\texttt{abort} & Fail the session. \\ \hline
\end{longtable}
}

\msubsection{struct}{Parameter}
A declared session parameter.

{\small
\begin{longtable}{|p{3.9cm}|p{4.4cm}|p{7.5cm}|}

\hline
\rowcolor{gray!33} Field & Type & Description \\ \hline
\endhead
\texttt{name} & \pref{Name} & Parameter name, referenced as \texttt{\$\{param.name\}}. \\ \hline
\texttt{type} & \pref{String} & Informative type: \texttt{string}, \texttt{number}, \texttt{boolean} or \texttt{object}. Optional. \\ \hline
\texttt{required} & \pref{Boolean} & Whether \texttt{workflowInstance.start} refuses a session without it. Optional. \\ \hline
\end{longtable}
}

\msubsection{struct}{PlanAction}
\texttt{ActionRequestDTO} of the Choreographer.

{\small
\begin{longtable}{|p{3.9cm}|p{4.4cm}|p{7.5cm}|}

\hline
\rowcolor{gray!33} Field & Type & Description \\ \hline
\endhead
\texttt{name} & \pref{Name} & Action name. \\ \hline
\texttt{nextActionName} & \pref{Name} & Next action. Optional. \\ \hline
\texttt{firstStepNames} & \pref{List}$<$\pref{Name}$>$ & Steps without predecessors. \\ \hline
\texttt{steps} & \pref{List}$<$\mref{PlanStep}$>$ & Steps. \\ \hline
\end{longtable}
}

\msubsection{struct}{PlanServiceRequirement}
Subset of the orchestration service requirement.

{\small
\begin{longtable}{|p{3.9cm}|p{4.4cm}|p{7.5cm}|}

\hline
\rowcolor{gray!33} Field & Type & Description \\ \hline
\endhead
\texttt{serviceDefinition} & \pref{ServiceDefinition} & Service. \\ \hline
\texttt{versions} & \pref{List}$<$\pref{String}$>$ & Accepted versions. Optional. \\ \hline
\texttt{interfaceTemplateNames} & \pref{List}$<$\pref{String}$>$ & Accepted interface templates. Optional. \\ \hline
\end{longtable}
}

\msubsection{struct}{PlanStep}
\texttt{StepRequestDTO} of the Choreographer.

{\small
\begin{longtable}{|p{3.9cm}|p{4.4cm}|p{7.5cm}|}

\hline
\rowcolor{gray!33} Field & Type & Description \\ \hline
\endhead
\texttt{name} & \pref{Name} & Step name. \\ \hline
\texttt{quantity} & \pref{Integer} & Always 1. Optional. \\ \hline
\texttt{serviceRequirement} & \mref{PlanServiceRequirement} & Service of the step. \\ \hline
\texttt{localCloudOnly} & \pref{Boolean} & From the step policy. Optional. \\ \hline
\texttt{staticParameters} & \pref{Map}$<$\pref{String}, \pref{String}$>$ & The step's \texttt{params} as strings. Optional. \\ \hline
\texttt{nextStepNames} & \pref{List}$<$\pref{Name}$>$ & Successors. Optional. \\ \hline
\end{longtable}
}

\msubsection{struct}{Policy}
Execution policy of a step. Every field other than \texttt{lease} and \texttt{localCloudOnly} is an extension outside the proven model.

{\small
\begin{longtable}{|p{3.9cm}|p{4.4cm}|p{7.5cm}|}

\hline
\rowcolor{gray!33} Field & Type & Description \\ \hline
\endhead
\texttt{quantity} & \pref{Integer} & Number of executions; only 1 is accepted (check M01). Optional. \\ \hline
\texttt{timeoutSeconds} & \pref{Duration} & Extension x-timeout: the step fails when it runs longer. Optional. \\ \hline
\texttt{nominalDurationSeconds} & \pref{Duration} & Expected duration, used to check lease durations (check L04). Optional. \\ \hline
\texttt{retry} & \mref{Retry} & Extension x-rebind: behaviour on \texttt{ERROR}. Optional. \\ \hline
\texttt{lease} & \mref{LeasePolicy} & The step's role for a shared resource. Optional. \\ \hline
\texttt{localCloudOnly} & \pref{Boolean} & Whether orchestration may return providers of other clouds (default true). Optional. \\ \hline
\end{longtable}
}

\msubsection{struct}{Recipe}
A workflow recipe $W$~\cite{varga2026}: named, versioned and immutable once registered. Actions run in sequence (\texttt{firstAction}, then each action's \texttt{next}); the steps of one action form a directed acyclic graph.

{\small
\begin{longtable}{|p{3.9cm}|p{4.4cm}|p{7.5cm}|}

\hline
\rowcolor{gray!33} Field & Type & Description \\ \hline
\endhead
\texttt{name} & \pref{Name} & Recipe name. \\ \hline
\texttt{version} & \pref{Version} & Recipe version; a registered name and version never changes content. \\ \hline
\texttt{description} & \pref{String} & Free text. Optional. \\ \hline
\texttt{firstAction} & \pref{Name} & The action $\alpha_0$ that is active when a session starts. \\ \hline
\texttt{parameters} & \pref{List}$<$\mref{Parameter}$>$ & Declared session parameters $\rho$. Optional. \\ \hline
\texttt{resources} & \pref{List}$<$\mref{Resource}$>$ & Shared resources (tools) used under leases. Optional. \\ \hline
\texttt{dispatch} & \mref{DispatchMode} & Default dispatch mode of the recipe; overridden per action or per session. Optional. \\ \hline
\texttt{actions} & \pref{List}$<$\mref{Action}$>$ & The actions $A$ with their steps. \\ \hline
\end{longtable}
}

\msubsection{struct}{RecipeFilter}
Optional filter of the recipe query.

{\small
\begin{longtable}{|p{3.9cm}|p{4.4cm}|p{7.5cm}|}

\hline
\rowcolor{gray!33} Field & Type & Description \\ \hline
\endhead
\texttt{name} & \pref{Name} & Only versions of this recipe. Optional. \\ \hline
\end{longtable}
}

\msubsection{struct}{RecipeKey}
Identifies one registered recipe version.

{\small
\begin{longtable}{|p{3.9cm}|p{4.4cm}|p{7.5cm}|}

\hline
\rowcolor{gray!33} Field & Type & Description \\ \hline
\endhead
\texttt{name} & \pref{Name} & Recipe name. \\ \hline
\texttt{version} & \pref{Version} & Recipe version; for \texttt{compile} the latest active version when absent. \\ \hline
\end{longtable}
}

\msubsection{struct}{RecipeList}
Registered recipes in registration order.

{\small
\begin{longtable}{|p{3.9cm}|p{4.4cm}|p{7.5cm}|}

\hline
\rowcolor{gray!33} Field & Type & Description \\ \hline
\endhead
\texttt{entries} & \pref{List}$<$\mref{RecipeRecord}$>$ & The recipes. \\ \hline
\texttt{count} & \pref{Integer} & Number of entries. \\ \hline
\end{longtable}
}

\msubsection{struct}{RecipeRecord}
A registered recipe as stored by the WorkflowManager.

{\small
\begin{longtable}{|p{3.9cm}|p{4.4cm}|p{7.5cm}|}

\hline
\rowcolor{gray!33} Field & Type & Description \\ \hline
\endhead
\texttt{id} & \pref{Integer} & Database identifier. \\ \hline
\texttt{name} & \pref{Name} & Recipe name. \\ \hline
\texttt{version} & \pref{Version} & Recipe version. \\ \hline
\texttt{status} & \mref{RecipeStatus} & \texttt{ACTIVE} or \texttt{RETIRED}. \\ \hline
\texttt{registeredBy} & \pref{SystemName} & The system that registered the recipe. \\ \hline
\texttt{recipe} & \mref{Recipe} & The canonical recipe document. \\ \hline
\texttt{createdAt} & \pref{DateTime} & Registration time. \\ \hline
\end{longtable}
}

\msubsection{struct}{RecipeRegistration}
Result of a registration.

{\small
\begin{longtable}{|p{3.9cm}|p{4.4cm}|p{7.5cm}|}

\hline
\rowcolor{gray!33} Field & Type & Description \\ \hline
\endhead
\texttt{recipe} & \mref{RecipeRecord} & The stored recipe; null when the recipe was not accepted. Optional. \\ \hline
\texttt{validation} & \mref{ValidationReport} & The validation report. \\ \hline
\texttt{created} & \pref{Boolean} & True when a new version was stored; false when identical content already existed. \\ \hline
\end{longtable}
}

\msubsection{enum}{RecipeStatus}
Life-cycle state of a registered recipe.

{\small
\begin{longtable}{|p{3.9cm}|p{11.9cm}|}

\hline
\rowcolor{gray!33} Value & Meaning \\ \hline
\endhead
\texttt{ACTIVE} & Sessions may be started. \\ \hline
\texttt{RETIRED} & Kept for traceability; new sessions are refused with 409. \\ \hline
\end{longtable}
}

\msubsection{struct}{Resource}
A shared resource that steps reserve, use and release under a lease (Proposition~1 of~\cite{varga2026}, lease variant).

{\small
\begin{longtable}{|p{3.9cm}|p{4.4cm}|p{7.5cm}|}

\hline
\rowcolor{gray!33} Field & Type & Description \\ \hline
\endhead
\texttt{name} & \pref{Name} & Resource name. \\ \hline
\texttt{leaseDurationSeconds} & \pref{Duration} & Lease duration granted by \texttt{workflowLease}; must exceed the longest tool-using step. Optional. \\ \hline
\texttt{provider} & \pref{SystemName} & Lease provider; the WorkflowManager when absent. Optional. \\ \hline
\end{longtable}
}

\msubsection{struct}{Retry}
Extension x-rebind.

{\small
\begin{longtable}{|p{3.9cm}|p{4.4cm}|p{7.5cm}|}

\hline
\rowcolor{gray!33} Field & Type & Description \\ \hline
\endhead
\texttt{maxAttempts} & \pref{Integer} & Maximum attempts of the step in one session. Optional. \\ \hline
\texttt{onError} & \mref{OnError} & What an \texttt{ERROR} report causes. Optional. \\ \hline
\end{longtable}
}

\msubsection{struct}{ServiceRequirement}
Requirements on the provider of the step's service, passed to orchestration.

{\small
\begin{longtable}{|p{3.9cm}|p{4.4cm}|p{7.5cm}|}

\hline
\rowcolor{gray!33} Field & Type & Description \\ \hline
\endhead
\texttt{serviceDefinition} & \pref{ServiceDefinition} & The service to consume. \\ \hline
\texttt{interfaces} & \pref{List}$<$\pref{String}$>$ & Accepted interface template names, e.g. \texttt{generic\_https}. Optional. \\ \hline
\texttt{minVersion} & \pref{String} & Lowest accepted service version. Optional. \\ \hline
\texttt{maxVersion} & \pref{String} & Highest accepted service version. Optional. \\ \hline
\texttt{metadata} & \pref{Map}$<$\pref{String}, \pref{JsonValue}$>$ & Metadata requirements on the service instance. Optional. \\ \hline
\texttt{qos} & \pref{Map}$<$\pref{String}, \pref{JsonValue}$>$ & QoS requirements forwarded to orchestration. Optional. \\ \hline
\texttt{requiresResults} & \pref{List}$<$\pref{Name}$>$ & Steps whose results the provider needs; each must be in $dep(s)$ (check D04). Optional. \\ \hline
\end{longtable}
}

\msubsection{struct}{Step}
A step $s$: one service exchange.

{\small
\begin{longtable}{|p{3.9cm}|p{4.4cm}|p{7.5cm}|}

\hline
\rowcolor{gray!33} Field & Type & Description \\ \hline
\endhead
\texttt{name} & \pref{Name} & Step name, unique in the recipe. \\ \hline
\texttt{service} & \mref{ServiceRequirement} & $\theta(s)$: the service the step consumes. \\ \hline
\texttt{next} & \pref{List}$<$\pref{Name}$>$ & Successors in the same action (the transitions $s \rightarrow s'$). Optional. \\ \hline
\texttt{dep} & \pref{List}$<$\pref{Name}$>$ & $dep(s)$: steps whose results the step reads. Each must run before $s$. Optional. \\ \hline
\texttt{executor} & \mref{ExecutorRequirement} & $\eta(s)$: which executor may run the step; absent means any executor declaring the service. Optional. \\ \hline
\texttt{params} & \pref{Map}$<$\pref{String}, \pref{JsonValue}$>$ & Static parameters passed unchanged to the provider. Optional. \\ \hline
\texttt{input} & \pref{Template} & $\kappa_s$: the request body sent to the provider, composed from $\rho$ and $\xi|dep(s)$. Optional. \\ \hline
\texttt{policy} & \mref{Policy} & Execution policy (quantity, timeout, retry, lease, local cloud). Optional. \\ \hline
\texttt{guards} & \pref{Map}$<$\pref{Name}, \pref{Expression}$>$ & Conditions on the transitions to successors (extension x-guard). A successor without a guard is always taken. Optional. \\ \hline
\end{longtable}
}

\msubsection{struct}{ValidationReport}
The result of validating a recipe. The verdict is the worst outcome of the findings (REJECT $>$ CLARIFY $>$ REVISE $>$ ACCEPT).

{\small
\begin{longtable}{|p{3.9cm}|p{4.4cm}|p{7.5cm}|}

\hline
\rowcolor{gray!33} Field & Type & Description \\ \hline
\endhead
\texttt{recipe} & \pref{String} & \texttt{name@version} of the validated recipe. \\ \hline
\texttt{verdict} & \mref{Verdict} & Overall verdict; only \texttt{ACCEPT} can be registered. \\ \hline
\texttt{findings} & \pref{List}$<$\mref{Finding}$>$ & Individual findings; empty when the recipe has none. \\ \hline
\end{longtable}
}

\msubsection{enum}{Verdict}
Validator outcome.

{\small
\begin{longtable}{|p{3.9cm}|p{11.9cm}|}

\hline
\rowcolor{gray!33} Value & Meaning \\ \hline
\endhead
\texttt{ACCEPT} & The recipe can be registered (also used for informational findings). \\ \hline
\texttt{REVISE} & The recipe must be changed. \\ \hline
\texttt{CLARIFY} & The recipe refers to something that cannot be confirmed (undeclared parameter, unknown service or executor). \\ \hline
\texttt{REJECT} & The recipe cannot be used at all. \\ \hline
\end{longtable}
}


\subsection{Primitives}
\label{sec:model:primitives}
Types and structures mentioned throughout this document that are assumed to be available to implementations of this
service. The concrete interpretation of each of these types must be provided by any IDD document claiming to implement
this service.

{\small
\begin{longtable}{|p{3.2cm}|p{12.6cm}|}

\hline
\rowcolor{gray!33} Type & Description \\ \hline
\endhead
\pdef{Boolean} & One out of \texttt{true} or \texttt{false}. \\ \hline
\pdef{DateTime} & Pinpoints a moment in time (UTC). \\ \hline
\pdef{Duration} & A length of time in seconds. \\ \hline
\pdef{Expression} & A side-effect free condition over the session parameters and stored results (Section on guard expressions). \\ \hline
\pdef{Integer} & A whole number. \\ \hline
\pdef{JsonValue} & Any value: object, array, string, number, Boolean or null. \\ \hline
\pdef{List} & An ordered collection of items of the type given in angle brackets. \\ \hline
\pdef{Map} & A collection of key/value pairs of the types given in angle brackets. \\ \hline
\pdef{Name} & An identifier of a recipe, action, step, parameter or resource: a letter followed by at most 62 letters, digits, \texttt{-} or \texttt{\_}. \\ \hline
\pdef{ServiceDefinition} & The name of an Arrowhead service definition. Names are compared case-insensitively and without \texttt{-}, \texttt{\_} and spaces, so \texttt{temperature-read} and \texttt{temperatureRead} denote the same service. \\ \hline
\pdef{String} & A chain of characters. \\ \hline
\pdef{SystemName} & The name of an Arrowhead system as it appears in its certificate and in the ServiceRegistry. \\ \hline
\pdef{Template} & A value in which strings of the form \texttt{\$\{param.$p$\}} or \texttt{\$\{result.$s$.path\}} are replaced by a session parameter or by (a part of) a stored step result. \\ \hline
\pdef{Version} & A recipe version in the form MAJOR.MINOR.PATCH. \\ \hline
\end{longtable}
}


\newpage
\bibliographystyle{IEEEtran}
\bibliography{bibliography}

\newpage

\section{Revision History}
\subsection{Amendments}

\noindent\begin{tabularx}{\textwidth}{| p{1cm} | p{2.2cm} | p{1.6cm} | X | p{3.6cm} |} \hline
\rowcolor{gray!33} No. & Date & Version & Subject of Amendments & Author \\ \hline
1 & 2026-11-23 & \arrowversion & Initial edition: service description of workflowRecipe, written from the Java reference implementation (ah-Workflow-java). & Documentation maintenance \\ \hline
\end{tabularx}

\subsection{Quality Assurance}

\noindent\begin{tabularx}{\textwidth}{| p{1cm} | p{2.2cm} | p{1.6cm} | X |} \hline
\rowcolor{gray!33} No. & Date & Version & Approved by \\ \hline
1 & & \arrowversion & Pending review \\ \hline
\end{tabularx}

\end{document}
